\documentclass{article}
\usepackage{amsmath, amssymb, amsthm}
\usepackage{hyperref}
\usepackage{geometry}
\geometry{margin=1in}

\title{Erd\H{o}s Problem \#1}
\date{Last updated: 2025-08-31}
\author{Source: erdosproblems.com}

\begin{document}
\maketitle

\noindent\textbf{Status:} OPEN \quad \textbf{Prize: \$500}

\noindent\textbf{Tags:} number theory, additive combinatorics

\medskip
\noindent\textbf{Problem Statement:}

If $A\subseteq \{1,\ldots,N\}$ with $\lvert A\rvert=n$ is such that the subset sums $\sum_{a\in S}a$ are distinct for all $S\subseteq A$ then\[N \gg 2^{n}.\]

\bigskip
\noindent\textbf{Remarks:}

Erd\H{o}s called this 'perhaps my first serious problem' (in \cite{Er98} he dates it to 1931). The powers of $2$ show that $2^n$ would be best possible here: this provides an upper bound for the minimal such $N$ of $N\leq 2^{n-1}$. This was improved by Conway and Guy \cite{CoGu68} (see also \cite{Gu82}) to $N\leq 2^{n-2}$ (for all large $n$). The best known upper bound is\[N\leq 0.22002\cdot 2^n,\]due to Bohman.

The trivial lower bound is $N \gg 2^{n}/n$, since all $2^n$ distinct subset sums must lie in $[0,Nn)$. Erd\H{o}s and Moser \cite{Er56} proved\[ N\geq (\tfrac{1}{4}-o(1))\frac{2^n}{\sqrt{n}}.\](In \cite{Er85c} Erd\H{o}s offered \$100 for any improvement of the constant $1/4$ here.)

A number of improvements of the constant have been given (see \cite{St23} for a history), with the current record $\sqrt{2/\pi}$ first proved in unpublished work of Elkies and Gleason. Two proofs achieving this constant are provided by Dubroff, Fox, and Xu \cite{DFX21}, who in fact prove the exact bound $N\geq \binom{n}{\lfloor n/2\rfloor}$.

An equivalent formulation is to ask for the maximal size of a set of integers in $[1,x]$ which is dissociated (all subset sums are distinct). If $F(x)$ is the size of such a set then this problem is equivalent to\[F(x) <\log_2x+O(1).\]Conway and Guy (see \cite{Gu82}) conjectured that $F(2^k)=k+2$ for all large $k$, but Erd\H{o}s \cite{Er80} wrote he had 'no opinion'.

In \cite{Er73} and \cite{ErGr80} the generalisation where $A\subseteq (0,N]$ is a set of real numbers such that the subset sums all differ by at least $1$ is proposed, with the same conjectured bound. (The second proof of \cite{DFX21} applies also to this generalisation.) This generalisation seems to have first appeared in \cite{Gr71}.

This problem appears in Erd\H{o}s' book with Spencer \cite{ErSp74} in the final chapter titled 'The kitchen sink'. As Ruzsa writes in \cite{Ru99} "it is a rich kitchen where such things go to the sink". 

The sequence of minimal $N$ for a given $n$ is A276661 in the OEIS.

See also [350].

This is discussed in problem C8 of Guy's collection \cite{Gu04}.

\bigskip
\noindent\textbf{References:}

[CoGu68] J. H. Conway and R. K. Guy, Sets of natural numbers with distinct sums. Notices Amer. Math. Soc. (1968), 345.
 
 [DFX21] Dubroff, Q. and Fox, J. and Xu, M. W., A note on the Erd\H{o}s distinct subset sums problem. SIAM Journal on Discrete Mathematics (2021), 322-324.
 
 [Er56] Erd\H{o}s, P., Problems and results in additive number theory. Colloque sur la Th\'{e}orie des Nombres, Bruxelles, 1955 (1956), 127-137.
 
 [Er73] Erd\H{o}s, P., Problems and results on combinatorial number theory. A survey of combinatorial theory (Proc. Internat. Sympos., Colorado State Univ., Fort Collins, Colo., 1971) (1973), 117-138.
 
 [Er80] Erd\H{o}s, Paul, A survey of problems in combinatorial number theory. Ann. Discrete Math. (1980), 89-115.
 
 [Er85c] Erd\H{o}s, P., On some of my problems in number theory I would most like to see solved. Number theory (Ootacamund, 1984) (1985), 74-84.
 
 [Er98] Erd\H{o}s, Paul, Some of my new and almost new problems and results in combinatorial number theory. Number theory (Eger, 1996) (1998), 169-180.
 
 [ErGr80] Erd\H{o}s, P. and Graham, R., Old and new problems and results in combinatorial number theory. Monographies de L'Enseignement Mathematique (1980).
 
 [ErSp74] Erd\H{o}s, Paul and Spencer, Joel, Probabilistic methods in combinatorics. Akad\'{e}miai Kiad\'{o} (1974).
 
 [Gr71] Graham, R. L., On sums of integers taken from a fixed sequence. (1971), 22--40.
 
 [Gu04] Guy, Richard K., Unsolved problems in number theory. (2004), xviii+437.
 
 [Gu82] Guy, Richard K., Sets of integers whose subsets have distinct sums. (1982), 141--154.
 
 [Ru99] Ruzsa, I., Erd\H{o}s and the Integers. Journal of Number Theory (1999), 115-163.
 
 [St23] Steinerberger, S., Some remarks on the Erd\H{o}s distinct subset sums problem. arXiv:2208.12182 (2023).

\bigskip
\noindent\small{Source: \url{https://www.erdosproblems.com/1}}
\end{document}
